\documentclass{leveldoc}
\DocUses{foundations,slingshot,first_coil,first_magnet}
\DocFingerprint{441a120cc2fa5805}

\begin{document}

\section{What the level asks}\label{sec:asks}

The game's text:
\leveldescription

Two neutral particles are launched from the left edge at $(\val{xZero}, \val{yZero})$ along
$+x$: the two spin states, each with a magnetic moment perpendicular to the board,
$\mu_\mathrm{p} = \val{momUp}$ (up, along $+z$) and $\mu_\mathrm{p} = \val{momDown}$ (down).
Without charge they feel no Lorentz force; only the gradient of the magnetic field acts on
them. Spin up must reach the lower detector, $\val{xDetLoUp} \le x \le \val{xDetHiUp}$,
$\val{yDetLoUp} \le y \le \val{yDetHiUp}$; spin down the upper one,
$\val{xDetLoDown} \le x \le \val{xDetHiDown}$, $\val{yDetLoDown} \le y \le \val{yDetHiDown}$
(no acceptance windows, no kinetic windows, no gates). The player places up to three magnets of
the listed strengths, of either orientation, in the region $\val{regX0} \le x \le \val{regX1}$,
$\val{regY0} \le y \le \val{regY1}$. A flight ends when the particle enters its detector, touches
a magnet, or leaves the world (\xsecref{foundations}{sec:board}).

With nothing placed both fly straight out through the right edge. The reference places one
magnet ($\odot$) above the launch line: its field is strongest near it, so the spin whose energy
falls in a stronger field is drawn up into the upper detector, and the other is pushed down into
the lower one. This document derives the force on a moment, writes it as a central potential
about the magnet, and finds each flight twice: by integrating the equations of motion
numerically, and by reducing them to quadratures with the energy and the angular momentum.

\section{Setup}\label{sec:setup}

The level file, field by field (the appendix reproduces it):
\setuptable

Both particles have mass $m = \val{m}$ and charge $0$; $c = \val{c}$; each is launched with
$T_0 = \val{Tzero}$, so $\gamma_0 = \val{gammaZero}$, $p_0 = \val{pZero}$ and
$v_0 = \val{vZero}$ (\xeqref{foundations}{eq:launch}). The reference magnet stands at
$\vect{r}_\mathrm{m} = (\val{xMag}, \val{yMag})$ with the strength $\mu = \val{muMag}$ (field
units, \physref{2.2}) and the radius $\val{Rm}$. We write $k = \mu_\mathrm{p}\mu$: $k = \val{kUp}$
for spin up and $\val{kDown}$ for spin down. The world is $-1 \le x \le \val{xWorld}$,
$-1 \le y \le \val{yWorld}$.

\begin{figure}[htbp]
\centering
\begin{tikzpicture}
\begin{axis}[axis equal image, width=0.92\linewidth, xmin=-1.6, xmax=31.6, ymin=-1.6, ymax=21.6,
  xtick={0,5,...,30}, ytick={0,5,...,20}, tick label style={font=\small},
  xlabel={$x$}, ylabel={$y$}, label style={font=\small}]
  \draw[black!8] (0,0) grid (30,20);
  \draw[black!45, dashed] (-1,-1) rectangle (\rawval{xWorld},\rawval{yWorld});
  \draw[black!30, dotted] (\rawval{regX0},\rawval{regY0}) rectangle (\rawval{regX1},\rawval{regY1});
  \fill[green!50!black, fill opacity=0.18] (\rawval{xDetLoUp},\rawval{yDetLoUp}) rectangle (\rawval{xDetHiUp},\rawval{yDetHiUp});
  \draw[green!40!black] (\rawval{xDetLoUp},\rawval{yDetLoUp}) rectangle (\rawval{xDetHiUp},\rawval{yDetHiUp});
  \fill[green!50!black, fill opacity=0.18] (\rawval{xDetLoDown},\rawval{yDetLoDown}) rectangle (\rawval{xDetHiDown},\rawval{yDetHiDown});
  \draw[green!40!black] (\rawval{xDetLoDown},\rawval{yDetLoDown}) rectangle (\rawval{xDetHiDown},\rawval{yDetHiDown});
  \draw[black!40, dashed] (\rawval{xMag},\rawval{yMag}) circle[radius=\rawval{rTUp}];
  \draw[black!40, densely dotted] (\rawval{xMag},\rawval{yMag}) circle[radius=\rawval{rTDown}];
  \addplot[very thick, blue!70!black] table {pathup.dat};
  \addplot[very thick, red!70!black] table {pathdown.dat};
  \fill[black!70] (\rawval{xMag},\rawval{yMag}) circle[radius=\rawval{Rm}];
  \fill (\rawval{xZero},\rawval{yZero}) circle[radius=0.15];
\end{axis}
\end{tikzpicture}
\caption{The magnet (black), the two detectors (green), and the two flights: spin up (blue),
pushed away from the magnet into the lower detector, and spin down (red), drawn towards it into
the upper one. The dashed and dotted circles have the radii of the two turning points
\eqref{eq:turning}; each flight touches its circle.}
\label{fig:board}
\end{figure}

\section{Principles and model}\label{sec:model}

\begin{principles}
  \item A magnetic moment's energy in a field, $U = -\vect{\mu}_\mathrm{p}\cdot\vect{B}$, and the
    force $\vect{F} = \nabla(\vect{\mu}_\mathrm{p}\cdot\vect{B})$ on a moment of fixed direction
    (\Jackson{§5.7}).
  \item The magnet's field in the plane, $B_z = -\mu/r^3$ (\xeqref{first_magnet}{eq:dipole-plane}),
    perpendicular to the board.
  \item The relativistic equation of motion \xeqref{foundations}{eq:eom} with that force, exact
    for a moment perpendicular to the plane in a magnetic field along $z$ (\physref{3.2}).
  \item The board's rules; the classical Runge--Kutta method, cubic Hermite interpolation and the
    error estimate by doubling the step (\xsecref{foundations}{sec:board},
    \xsecref{foundations}{sec:numerics}), as in \xsecref{slingshot}{sec:method}.
\end{principles}

\begin{assumption}{Two spin states as two classical moments}\label{as:spin}
Each state is a classical moment of fixed direction, $\pm\uvec{z}$. In the plane the field is
along $z$, parallel to the moment, so there is no torque and no precession: the state stays a
state. The quantum superposition that splits one beam into two is replaced by two shots. The
game does the same (\physref{3.2}).
\end{assumption}

\begin{assumption}{No radiation}\label{as:radiation}
A moving moment radiates as an electric dipole $\vect{v}\times\vect{\mu}_\mathrm{p}/c^2$; the
power, of order $\mu_\mathrm{p}^2\abs{\dot{\vect{a}}}^2/c^7$, is far below any effect on the
flight. The game estimates it along every flight and checks it below $10^{-10}$ of $T_0$
(\physref{3.2}).
\end{assumption}

\section{Derivation}\label{sec:derivation}

\subsection{The force on a moment, as a potential}\label{sec:potential}

A moment $\mu_\mathrm{p}\uvec{z}$ in the field $B_z\uvec{z}$ has the energy
$U = -\mu_\mathrm{p}B_z$. Its interaction with a moving particle is the Lagrangian term
$\vect{\mu}_\mathrm{p}\cdot\vect{B}'/\gamma$, with $\vect{B}'$ the field in the particle's rest
frame; for a magnetic field along $z$, no electric field and motion in the plane,
$B'_z = \gamma B_z$, and the term is $\mu_\mathrm{p}B_z$, independent of the velocity
(\physref{3.2}). So the motion is that of a relativistic particle in the static potential
$U$:
\begin{equation}\label{eq:potential}
  U(\vect{x}) = -\mu_\mathrm{p}B_z = \frac{\mu_\mathrm{p}\mu}{r^3} = \frac{k}{r^3},
  \qquad r = \abs{\vect{x} - \vect{r}_\mathrm{m}},
\end{equation}
\begin{equation}\label{eq:eom-moment}
  \deriv{\vect{x}}{t} = \frac{\vect{p}}{\gamma m}, \qquad
  \deriv{\vect{p}}{t} = -\nabla U = \frac{3k\,(\vect{x} - \vect{r}_\mathrm{m})}{r^5} ,
\end{equation}
summed over the magnets. For $k > 0$ (spin up, with the magnet $\odot$) $U$ is positive and
falls with $r$: the particle is pushed away, towards weaker $\abs{B_z}$. For $k < 0$ it is drawn
in, towards stronger field. The force falls as $1/r^4$: the effect is a gradient's, strong only
close to the magnet. At the launch $B_z = \val{bZero}$ and $U = \val{UZeroUp}$ for spin up.

\subsection{Two conserved quantities}\label{sec:conserved}

$U$ does not depend on time, so the energy is conserved,
\begin{equation}\label{eq:energy-moment}
  W = \gamma mc^2 + \frac{k}{r^3} = \sqrt{m^2c^4 + p^2c^2} + \frac{k}{r^3},
\end{equation}
and $U$ depends only on $r$, so the force is central and the angular momentum about the magnet,
\begin{equation}\label{eq:angular}
  L = \big((\vect{x} - \vect{r}_\mathrm{m})\times\vect{p}\big)_z ,
\end{equation}
is conserved too. At the launch $W = \val{WUp}$ (up) and $\val{WDown}$ (down), and $L = \val{LUp}$
for both.

\subsection{The flights by quadratures}\label{sec:quadrature}

At the radius $r$ energy gives the momentum's size, $p(r)^2 = ((W - k/r^3)/c)^2 - m^2c^2$, and
angular momentum its azimuthal part, $p_\varphi = L/r$, so
\begin{equation}\label{eq:pr}
  p_r^2 = \Big(\frac{W - k/r^3}{c}\Big)^2 - m^2c^2 - \frac{L^2}{r^2}, \qquad
  \gamma(r) = \frac{W - k/r^3}{mc^2}.
\end{equation}
Both particles start moving inwards ($p_r = \val{prZeroUp}$) and turn where $p_r = 0$:
\begin{equation}\label{eq:turning}
  r_\mathrm{t}: \quad p_r^2(r_\mathrm{t}) = 0, \qquad
  r_\mathrm{t} = \val{rTUp} \text{ (up)}, \quad \val{rTDown} \text{ (down)},
\end{equation}
the largest root below the launch radius, found by bisection. (For spin down $p_r^2$ turns
positive again close to the magnet, where the attraction would win; the particle never gets
there.) From $\dot\varphi = L/(\gamma m r^2)$ and $\dot r = p_r/(\gamma m)$, on each branch
\begin{equation}\label{eq:phi-of-r}
  \deriv{\varphi}{r} = \frac{L}{r^2\,\abs{p_r}} ,
\end{equation}
\begin{equation}\label{eq:t-of-r}
  \deriv{t}{r} = \frac{\gamma(r)\,m}{\abs{p_r}} ,
\end{equation}
with $\dd r < 0$ on the way in and $\dd r > 0$ on the way out, exactly as in
\xsecref{first_magnet}{sec:quadrature}: the angle and the time on the outgoing branch are the
integrals from $r_\mathrm{t}$ to the launch radius plus those from $r_\mathrm{t}$ to $r$, made
smooth by $r = r_\mathrm{t} + u^2$ and evaluated to thirty digits; the entry into the detector is
the first point of the outgoing branch inside it, found by bisection on $r$. Spin up enters at
$r = \val{rEUp}$, at $t = \val{tQUp}$, at $y = \val{yQUp}$; spin down at $r = \val{rEDown}$, at
$t = \val{tQDown}$, at $y = \val{yQDown}$.

\subsection{The method}\label{sec:method}

The march integrates \eqref{eq:eom-moment} with the classical Runge--Kutta method
\xeqref{foundations}{eq:rk4} at the fixed step $h = \val{H}$, locates the events on the cubic Hermite
interpolant and estimates its error by doubling $h$ (\xsecref{slingshot}{sec:method}). $W$ and $L$
are its checks.

\section{The reference solution by hand}\label{sec:hand}

\begin{procedure}
  \item\label{st:launch} \emph{The launch.} $\gamma_0 = \val{gammaZero}$,
    $\vect{p}_0 = (\val{pZero}, 0)$, $r_0 = \val{rZero}$; $W$ by \eqref{eq:energy-moment}, $L$ by
    \eqref{eq:angular}.
  \item\label{st:turn} \emph{The turning points} \eqref{eq:turning}: $\val{rTUp}$ (up),
    $\val{rTDown}$ (down). The attracted spin comes closer.
  \item\label{st:up} \emph{Spin up.} After $\val{stepsUp}$ steps it enters its detector through
    the \val{faceUp} face at $t = \val{tUp}$, at $(\val{xUp}, \val{yUp})$, with momentum
    $(\val{pxUp}, \val{pyUp})$, $\val{angUp}$ degrees from $+x$. Step estimate $\val{errUp}$; $W$
    and $L$ drift by $\val{wDevUp}$ and $\val{LDevUp}$; the quadratures agree to
    $\val{quadDiffUp}$.
  \item\label{st:down} \emph{Spin down.} After $\val{stepsDown}$ steps it enters its detector
    through the \val{faceDown} face at $t = \val{tDown}$, at $(\val{xDown}, \val{yDown})$, with
    momentum $(\val{pxDown}, \val{pyDown})$, $\val{angDown}$ degrees from $+x$. Step estimate
    $\val{errDown}$; $W$ and $L$ drift by $\val{wDevDown}$ and $\val{LDevDown}$; the quadratures
    agree to $\val{quadDiffDown}$.
\end{procedure}

\begin{result}[The reference flights]
Spin up enters the lower detector at $(\val{xUp}, \val{yUp})$ at $t = \val{tUp}$; spin down
enters the upper detector at $(\val{xDown}, \val{yDown})$ at $t = \val{tDown}$. The level is
solved.
\end{result}

The split is not symmetric ($\val{angUp}$ against $\val{angDown}$ degrees): the attracted spin
passes closer to the magnet, where the $1/r^4$ force is stronger.

\section{Comparison with the game}\label{sec:comparison}

\comparisontable

Every difference is a few $10^{-12}$ or less, the size of the game's verified accuracy
(about $10^{-10}$, Foundations §6) and of our estimates (step $\val{errUp}$, quadratures
$\val{quadDiffUp}$). The tolerance, $10^{-9}$, is ten times the larger of the two computations'
estimates (or the game's own preview--verification gap where that is larger: here it is not),
fixed before comparing.

\section{Other solutions}\label{sec:others}

On the empty board, the second placement in \code{placements.json}, both spins fly straight
through the right edge at $(\val{bareX}, \val{bareY})$ at $t = \val{bareT}$: no field, no
gradient. One magnet near the path splits them; its orientation decides which spin goes which
way, and its strength and distance how far, roughly as $k/r_\mathrm{t}^4$ times the time spent
near it. A $\otimes$ magnet below the path would do what the reference's $\odot$ does above it.
A map of solutions would take the game's verdicts, which this document does not reproduce.

\section{What the game computes}\label{sec:engine}

The game gives a particle with a moment the force $\mu_\mathrm{p}\nabla B_z$, with the gradient of
the dipole field taken directly (\physref{3.2}), and integrates \eqref{eq:eom-moment} with its
adaptive eighth-order Runge--Kutta method (\physref{5}, \physref{6}, \physref{7}). It refuses
moments together with electric sources at finite $c$, where the rest-frame field would make the
interaction velocity dependent (\physref{3.2}); this level has none. Its conserved
$\gamma mc^2 - \mu_\mathrm{p}B_z$ and $L$ are \eqref{eq:energy-moment} and \eqref{eq:angular}.

\section{Further reading}\label{sec:reading}

The force and energy of a magnetic moment in a field: \Jackson{§5.7}. The magnet's field:
\levelref{first_magnet}; the quadratures for a central or axial symmetry:
\xsecref{first_magnet}{sec:quadrature}.

\levelfileappendix

\end{document}
